\documentclass[11pt]{article}
\usepackage{amsmath,amssymb,amsthm,booktabs,longtable,hyperref}
\newtheorem{theorem}{Theorem}
\newtheorem{lemma}[theorem]{Lemma}
\newtheorem{proposition}[theorem]{Proposition}
\newtheorem{corollary}[theorem]{Corollary}
\newcommand{\sq}{\sqcup}
\title{Nonisolated middle vertices in quintic graphs of diameter eight}
\author{Human Compression project}
\date{}
\begin{document}
\maketitle

\paragraph{Scope and attribution.}
We give a finite description relative to a fixed geodesic of length eight.
It is a description of marked incidence objects, not a count of ordinary
graphs. Different geodesics and different rows below may describe the same
ordinary graph. No graph enumeration, Source execution, TeX compilation,
or Lean verification was performed in preparing this note.

The original order--girth--diameter problem is due to Cambie, Goedgebeur,
Jooken and Van den Eede~\cite{cages}. The independent-neighborhood reduction,
the empty and singleton middle cases, and the six forms of one small side
are retained project results~\cite{middle}. Incidence margins, vertex
covers, graph distance, and omission multigraphs are standard ingredients.
The present contribution is their assembly into the finite nonisolated
middle list and the radius argument. No historical priority is claimed.

\section{Notation and local restrictions}

Let $G$ be simple, quintic, of order $32$, girth four and diameter eight.
Fix an eight-geodesic with endpoints $p_0,p_8$ and midpoint $p_4$, and put
\[
 A=N(p_0),\qquad B=N(p_4),\qquad C=N(p_8).
\]
We import the retained reduction: $I=A\sq B\sq C$ is independent of order
$15$; $J=V(G)\setminus I$ induces exactly five edges; the three displayed
vertices have no neighbors in $J$. The remaining fourteen vertices split
as $L,R,M$, where $L$ meets $A$, $R$ meets $C$, and $M$ meets neither.
Every vertex of $L$ or $R$ has positive incidence with its corresponding
endpoint rows. Moreover $|L|,|R|\geq5$. We write $F=G[J]$, omitting its
three named isolated vertices when discussing its shape.

The empty middle case is impossible, and the singleton middle case has
already been settled~\cite{middle}. Here $q=|M|\in\{2,3,4\}$ and at least
one vertex of $M$ is nonisolated in $F$. We assume $G$ is nonbipartite.
For an endpoint side $S$, let $D_S=F[S]$, let $e_S=|E(D_S)|$, and let
$x_S$ count its $F$ edges to $M$. An endpoint row is an independent
four-subset of its side; its fifth neighbor is $p_0$ or $p_8$.

\begin{lemma}\label{distance}
The distance between $L$ and $R$ is at least four. A vertex of $M$ cannot
have $F$ neighbors on both sides. A vertex of $B$ cannot meet both sides.
Partition $B$ into $B_L,B_R,B_0$ according as it meets $L$, meets $R$, or
meets neither. Both $B_L,B_R$ are nonempty. If $M_S$ consists of the middle
vertices with an $F$ neighbor on side $S$, then $B_L$ avoids $M_R$ and
$B_R$ avoids $M_L$. No $F$ edge joins $M_L$ to $M_R$.

A vertex of $B_0$ has four neighbors in $M$. Consequently $B_0$ is empty
unless $q=4$, and when it is nonempty every one of its vertices meets all
four middle vertices and $F[M]$ has no edges.
\end{lemma}
\begin{proof}
Every side vertex is at distance two from its endpoint. A path of at most
three edges between the sides would contradict $d(p_0,p_8)=8$. Each
restriction follows by writing the resulting path of length at most
three between the sides. The original geodesic supplies at least one
port on each side. Every $B$ vertex already meets $p_4$ and has exactly
four neighbors among the fourteen vertices. With no side neighbor these
must all lie in $M$. A common neighbor in $B_0$ would make a triangle
with any edge of $F[M]$.
\end{proof}

Conversely these listed separation restrictions exclude every path of
length at most three between the sides. This is checked by its possible
interior vertices: a $B$ vertex cannot join the sides; a middle vertex
cannot meet both sides; the only remaining three-edge possibilities
are $S$--$M_S$--$B_T$--$T$, $S$--$B_S$--$M_T$--$T$, or
$S$--$M_S$--$M_T$--$T$. Endpoint rows and the three named vertices add
no shorter possibility. A path through $B_S,p_4,B_T$ has length four
between suitable side vertices. Thus these local restrictions give
exactly the required endpoint distance once the row incidences exist.

\begin{lemma}\label{star}
In this setup $G$ is bipartite if and only if the five edges of $F$ form
a star $K_{1,5}$. In the nonbipartite case every $F$ degree is at most four.
\end{lemma}
\begin{proof}
In a bipartite quintic graph of order $32$ the parts have order $16$.
The vertices of $I$ lie in the same part, because the three named
vertices are at even distances on the fixed geodesic. The one remaining
vertex in that part has its five edges in $F$, giving the stated star.
Conversely the centre of such a star has no neighbor in $I$; moving it
from $J$ to $I$ gives a bipartition. A vertex of degree five in $F$
would use all five edges, so gives exactly that star.
\end{proof}

Write $b_S=|B_S|$ and $j=|B_0|$. A middle vertex outside $M_L\cup M_R$
will be called neutral; this describes its side incidences, not its
distance from the midpoint.

\begin{lemma}\label{edge}
For $uv\in E(F[M])$ one has $d_F(u)+d_F(v)\geq5$. A neutral middle
vertex has $F$ degree zero or one. If it has degree one, its neighbor
has degree at least four.
\end{lemma}
\begin{proof}
The $B$ neighborhoods of the adjacent vertices are disjoint. Their
sizes are $5-d_F(u)$ and $5-d_F(v)$, proving the first assertion.
A neutral vertex has at most $q-1\leq3$ neighbors, all in $M$. Suppose
its degree is $k=2$ or $3$. Each neighbor has degree at least $5-k$.
The neighbors are independent since $F$ has no triangle. Their incident
edge sets are therefore disjoint, and account for at least
$k(5-k)=6$ edges of $F$, a contradiction. Degree one gives the last claim.
\end{proof}

We shall also use the side-edge inequality. If $uv\in E(D_S)$, its two
endpoint-row incidences have total at most five and its two $B_S$
incidences have total at most $b_S$. Thus
\begin{equation}\label{sideedge}
 d_F(u)+d_F(v)\geq5-b_S.
\end{equation}
This uses the full degrees in $F$, including edges to $M$.

\section{When two middle vertices are adjacent}

\begin{proposition}\label{MM}
If $F[M]$ has an edge, then $F=K_{1,4}\sq K_2$, padded with isolates.
Up to reversal of the two endpoints its incidence data are exactly the
following seven rows. Here $u\in M_S$ is the degree-four star centre;
$r$ of its leaves are neutral middle vertices; the other $4-r$ leaves
lie in $S$; and $z$ further middle vertices are isolated in $F$.
\[
\begin{array}{c|c|c|c|c|c|c}
q&(|S|,|T|)&r&z&(b_S,b_T)&\text{remaining edge}&\text{range}\\\hline
3&(5,6)&r&2-r&(1,4)&D_T\text{ edge}&r=1,2\\
4&(5,5)&r&3-r&(2,3)&D_T\text{ edge}&r=1,2,3\\
4&(5,5)&r&2-r&(1,4)&w\!\in\!M_T\text{ to }T&r=1,2
\end{array}
\]
There is no such nonbipartite configuration for $q=2$.
\end{proposition}
\begin{proof}
Here $j=0$ by Lemma~\ref{distance}. Adjacent middle vertices cannot both
belong to $M_S$. Their disjoint $B$ neighborhoods would then lie in
$B_S$, so their degree sum would be at least $10-b_S$. On the other hand
an edge's endpoint degrees have sum at most six, since $F$ has five
edges. Hence $b_S=4$, $b_T=1$, the degree sum is six, and every edge
meets one of these two vertices. Side $T$ would have no $F$ incidences,
requiring $5|T|-20\geq5$ side-port incidences from just one port, which
has at most four. This is impossible. Opposite side types are excluded
by Lemma~\ref{distance}.

Thus every middle edge meets a neutral vertex. Lemmas~\ref{edge} and
\ref{star} show that this vertex is a leaf and the other endpoint has
degree four. All middle edges therefore belong to the same four-edge
star: a second degree-four centre outside it cannot be supported by the
one remaining edge. Let $a$ be the unique $B$ neighbor of its centre
$u\in M_S$. Every neutral leaf meets exactly $B\setminus\{a\}$.
In particular $a\in B_S$. The remaining edge cannot meet $u$, a middle
leaf, or a third middle vertex already in the star. It is either a
side edge, or an edge from one further middle vertex $w$ to a side.
All other middle vertices are isolated and meet every $B$ vertex.

If the remaining edge is a side edge, put $z=q-1-r$. The port $a$ has
side arity $3-z$; every other port has arity $5-q$. If the edge lies
in $S$, side $T$ has no $F$ incidence, giving
\[
 b_T(5-q)=5|T|-20.
\]
For $q=2,3,4$ and $|T|\geq5$ this has no solution with
$1\leq b_T\leq4$ and the possible side sizes. If the edge lies in $T$,
the equation is instead
\[
 b_T(5-q)=5|T|-22.
\]
For $q=2$ its only possible small-side solution is $|T|=5,b_T=1$;
that side edge violates~\eqref{sideedge}. For $q=3$ it forces
$|T|=6,b_T=4$, and for $q=4$ it forces $|T|=5,b_T=3$.
These give the first two rows and their stated $r,z$ values.

Finally let the remaining edge meet a new middle vertex $w$. Its degree
is one, so it needs four neighbors in its own $B$ part. If $w\in M_S$,
then $b_S=4,b_T=1$. Side $T$ would require
$5|T|-20=6-q$, which is impossible. If $w\in M_T$, then
$b_T=4,b_S=1$. Now $z=q-2-r$, and the arity of $a$ is $3-z$.
Side $S$ gives
\[
 5|S|-20-(4-r)=3-z.
\]
Only $q=4$, $|S|=5$ and $r+z=2$ are possible. This proves the third row.
All remaining edges are disjoint from the four-edge star, as asserted.
\end{proof}

The complete $B$ pattern in this proposition is as follows. The isolated
middle vertices meet every port; the $r$ neutral leaves meet all ports
except $a$; the star centre meets only $a$. In the third row $w$ meets
all four ports in $B_T$. In the first row $a$ meets all $r+1$ side
vertices outside the star, and every $B_T$ port meets an independent
two-subset of its six-vertex side. In the second row $a$ meets an
$r$-subset outside the star, the other $B_S$ port meets one side vertex,
and the three $B_T$ ports meet individual endpoints of the side edge.
In the third row $a$ meets every side vertex outside the star, and
each $B_T$ port meets a singleton avoiding the side neighbor of $w$.
These conditions, together with the endpoint-row conditions below, are
both necessary and sufficient; they do not assign ordinary class counts.

\section{No edges between middle vertices}

Now a neutral middle vertex is isolated in $F$. Put $z=|M\setminus
(M_L\sq M_R)|$, $m_S=|M_S|$, and let the positive degrees in $M_S$ be
the multiset $\mathbf d_S$. Its sum is $x_S$. Every isolated middle
vertex meets all five ports; every $B_0$ port meets all middle vertices;
and a vertex of degree $d$ in $M_S$ has exactly $5-j-d$ neighbors in
$B_S$. It has no neighbor in the opposite $B$ part. A $B_S$ port with
$k$ neighbors in $M_S$ has side arity $4-z-k$, which must be positive.
Thus
\begin{align}
 5-j-b_S&\leq d\leq\min(4,5-j),\label{degreebound}\\
 2(e_S+x_S)&=5|S|-20-b_S(4-z)+m_S(5-j).\label{balance}
\end{align}
The first bound also includes $d\geq1$. The second follows by counting
side-port incidences both from the five endpoint rows and from the ports.
Always $e_L+e_R+x_L+x_R=5$.

\begin{lemma}\label{five}
On a five-vertex side, $D_S$ has at most one edge. If it has an edge
$xy$, all other side vertices meet all five endpoint rows and have no
middle or port neighbors. Each vertex of $M_S$ has degree one and meets
$x$ or $y$; every $B_S$ port has side arity one and meets $x$ or $y$.
Consequently $b_S+m_S=3$.
\end{lemma}
\begin{proof}
Two disjoint side edges prevent an independent four-subset. Two adjacent
side edges force every independent four-subset to omit their common
vertex, violating its positive endpoint incidence. Hence there is at
most one edge. Every endpoint row must omit $x$ or $y$, so the other
three vertices occur in all five rows. Their remaining incidences are
zero. A middle vertex cannot meet both endpoints of the edge, and a
port cannot contain both. Counting the remaining column deficits gives
$2+m_S+b_S=5$.
\end{proof}

For $q=2,3$, a five-vertex side has no internal edge in the present
nonisolated case. Indeed a nonempty $M_S$ would require
$b_S\geq4$ for its degree-one vertices, contrary to $b_S+m_S=3$.
If $m_S=0$, all its ports have arity $4-z=1$, giving $z=3$ and no
nonisolated middle vertex at all.

\begin{proposition}\label{smallq}
For $q=2$, the no-middle-edge cases are the following two rows, up to
endpoint reversal. Both sides have order six, $z=1$, and all nonisolated
middle vertices consist of one vertex $u\in M_S$.
\[
\begin{array}{c|c|c|c|c}
d_F(u)&(b_S,b_T)&D_S&D_T&\text{port arities on }T\\\hline
2&(3,2)&K_2&P_3\sq3K_1&3,3\\
3&(3,2)&\varnothing&P_3\sq3K_1&3,3
\end{array}
\]
Every $B_S$ port meets the isolated middle vertex. In the degree-two
row all three also meet $u$ and have side arity two. In the degree-three
row exactly two meet $u$ and have arity two; the third has arity three.

For $q=3$, the no-middle-edge cases have side sizes $(5,6)$, $z=2$, and
one vertex $u$ of degree $d$ on the five-vertex side $S$:
\[
\begin{array}{c|c|c|c}
d&(b_S,b_T)&D_T&\text{port arities on }T\\\hline
4&(1,4)&K_2\sq4K_1&2,2,2,2\\
3&(2,3)&P_3\sq3K_1\text{ or }2K_2\sq2K_1&2,2,2\\
2&(3,2)&P_4\sq2K_1&2,2\\
1&(4,1)&C_4\sq2K_1&2
\end{array}
\]
Here $D_S$ is empty. Each $B_S$ port meets $u$ and both isolated middle
vertices, and therefore meets one side vertex outside $N_F(u)$.
\end{proposition}
\begin{proof}
We include the finite integer check to make the list reproducible without
an auxiliary computation. For $q=2$, first suppose $z=1$ and put the
nonisolated vertex on $S$. A five-vertex $S$ would give
$d=(10-3b_S)/2\geq5-b_S$, an impossibility. A five-vertex opposite
side would require $3b_T=5$. Thus both sides have order six, and
\eqref{balance} forces $b_S=3$, $e_S+d=3$, and $e_T=2$.
The bound $d\geq2$ leaves exactly $d=2,3$.
Inequality~\eqref{sideedge} excludes two disjoint edges on $T$, leaving
$P_3$.

If $z=0$ and the two middle vertices are on the same side, parity in
\eqref{balance} forces side sizes $(6,6)$. The only possible port counts
are $(3,2)$ or $(4,1)$. The former forces two degree-two middle vertices
and one edge on the opposite side, which violates~\eqref{sideedge}.
The latter forces two degree-one middle vertices and three opposite
side edges. Inequality~\eqref{sideedge} leaves only a three-leaf star;
its centre occurs in no independent four-subset on six vertices.
This is the retained six-vertex obstruction. If the middle vertices
are split, parity forces side sizes $(5,7)$; the five-vertex side gives
$d=5-2b_S\geq5-b_S$, again impossible.

For $q=3$ let $S$ have order five. Then $e_S=0$ and
\[
 2x_S=5-(4-z)b_S+5m_S,\qquad
 x_S\geq m_S(5-b_S).
\]
For $z=2$ the only solution is $m_S=1$, $m_T=0$ and
$x_S=5-b_S$; $m_S=0$ would require $2b_S=5$.
For $z=1$, the choices $m_S=0,1,2$ respectively give a noninteger
port count, $b_S\leq0$, or $b_S\geq5$. For $z=0$, the choices
$m_S=0,1,2,3$ give respectively a noninteger port count, $b_S\leq0$,
an odd right-hand side, or $b_S\geq5$. Thus $z=2$ is forced.
Put $b_T=d$; the opposite side has $e_T=b_S=5-d$ and $d$ ports of
arity two. The side-edge inequality now requires endpoint degree sum
at least $5-d$ on every internal edge.

For one edge this gives the first row. For two edges both $P_3$ and
$2K_2$ survive. For three edges the only triangle-free possibilities
are $P_4,K_{1,3},P_3\sq K_2,3K_2$, padded with isolates. The last two
violate the edge inequality, and the star centre has no independent
four-subset, leaving $P_4$. For four edges on six vertices the possibilities
are the three five-vertex trees, $C_4$, $P_4\sq K_2$,
$K_{1,3}\sq K_2$, and $2P_3$. The edge inequality eliminates all but
$K_{1,4}$ and $C_4$; the star centre again has no endpoint incidence.
This leaves precisely the table.
\end{proof}

\section{The four-middle case without middle edges}

Both sides now have order five. First suppose neither has an internal
edge. The positive middle degrees sum to five, and no degree is five.
Their partitions are
\[
 (4,1),\ (3,2),\ (3,1,1),\ (2,2,1),\ (2,1,1,1).
\]
Positive middle degrees occur on both sides: if all were on one side,
the other side would give $b_T(4-z)=5$, impossible for
$4-z\in\{2,3,4\}$. Applying~\eqref{balance} gives exactly the following
seven rows. The two degree lists refer to $M_S,M_T$; the isolated
middle vertices are not in those lists. The last two columns give, for
each middle vertex in the corresponding list, its number of neighbors
in its own $B$ part. Its other $j$ port neighbors are in $B_0$.
\[
\begin{array}{c|c|c|c|c|c|c|c}
 &z&j&\mathbf d_S&\mathbf d_T&(b_S,b_T)&\mathbf r_S&\mathbf r_T\\\hline
 Q_1&2&0&(4)&(1)&(1,4)&(1)&(4)\\
 Q_2&2&0&(3)&(2)&(2,3)&(2)&(3)\\
 Q_3&2&2&(3)&(2)&(1,2)&(0)&(1)\\
 Q_4&1&1&(3)&(1,1)&(1,3)&(1)&(3,3)\\
 Q_5&1&2&(1)&(2,2)&(2,1)&(2)&(1,1)\\
 Q_6&1&3&(2)&(2,1)&(1,1)&(0)&(0,1)\\
 Q_7&0&2&(2)&(1,1,1)&(1,2)&(1)&(2,2,2)
\end{array}
\]
These specify the entire middle--port incidence pattern up to permutations
within the displayed equal-role classes. In $Q_3$ the degree-two middle
vertex chooses one of two ports, and that choice is equivalent under
their interchange. All zero or full entries in the table are literal.

\begin{proposition}\label{q4}
The seven rows above, together with the following three rows, are the
complete $q=4$ list with $F[M]$ empty and some middle vertex nonisolated.
For the additional rows $j=1,2,3$ and
\[
\begin{array}{c|c|c|c|c|c}
z&(b_S,b_T)&\mathbf d_S&\mathbf d_T&D_S&D_T\\\hline
4-j&(4-j,1)&(1^{j-1})&(5-j)&K_2&\varnothing.
\end{array}
\]
Each middle vertex on $S$ meets all $B_S$ and all $B_0$; the middle
vertex on $T$ meets only $B_0$. For $j\geq2$, every $B_S$ port meets
one endpoint of $D_S$ and every $M_S$ vertex meets its other endpoint.
For $j=1$ the three singleton side incidences of $B_S$ can be split
between the two edge endpoints. The single $B_T$ port has side arity $j$.
\end{proposition}
\begin{proof}
With no internal side edge put $k=4-z$. For a side with $m_S$ positive
degrees summing to $x_S$, balance reads
\[
 b_S=\frac{5+m_S(5-j)-2x_S}{k}.
\]
For $k=2$ the split is $1+1$. Divisibility forces even $j$;
$j=0$ gives $Q_1,Q_2$, and $j=2$ leaves only $Q_3$ after the bounds
$0\leq5-j-d\leq b_S$. For $k=3$ orient the split as $1+2$, with
single degree $d$. Then $3b_S=10-j-2d$. For $j=0,1,2,3$,
divisibility requires respectively $d\equiv2,0,1,2\pmod3$.
The same degree bounds leave exactly $Q_4,Q_5,Q_6$; in the $j=2$
case the opposite degrees $(3,1)$ fail the degree-one bound, while
$(2,2)$ work. For $k=4$, the split $2+2$ has an odd numerator and is
impossible. In the split $1+3$, single degree two forces $j=2$ and
gives $Q_7$; single degree one forces $j=0$ and violates its port bound.
This proves the seven-row table.

Suppose now a side has an internal edge. By Lemma~\ref{five} it has
$b_S+m_S=3$ and every one of its middle vertices has degree one.
There cannot be an internal edge on both sides. Otherwise the total
five-edge count gives $z=1$, the two side equations give $j=2$, and
one side has $b_S=1,m_S=2$. Its middle vertices require two own-side
ports, a contradiction.

Assume only $S$ has an internal edge. If $m_S>0$, each of its middle
vertices needs $4-j$ own-side ports. Hence
$b_S=3-m_S\geq4-j$. Since $b_T=2-j+m_S\geq1$, equality is forced:
$m_S=j-1$, $b_S=4-j$, $b_T=1$. Every middle vertex on $S$ therefore
meets every $B_S$ port. Their side arity one gives $z=4-j$; one
middle vertex remains on $T$, with degree $5-j$. This gives $j=2,3$.
If $m_S=0$, its ports of arity one force $z=3$ and there is one
degree-four middle vertex on $T$. Its sole $B$ neighbor must belong
to $B_0$: an own-side port would already meet all four middle vertices
and have no side neighbor. Thus $j=1$, giving the remaining row.
The asserted edge endpoints are forced by triangle avoidance.
\end{proof}

\section{A complete incidence construction}

The tables specify finite core data, not complete ordinary isomorphism
classes. Here are all the remaining choices, and their exact tests.

For each middle vertex with positive side degree choose that many
distinct side neighbors. For each port choose its prescribed middle
neighbors and a side subset of the prescribed arity. Choose the small
internal side graph in its indicated isomorphism type. All chosen port
neighbor sets must be independent in $F$. In particular a port adjacent
to a middle vertex must avoid that middle vertex's side neighbors.
There are no side-crossing $F$ edges and no incidences to an opposite
side's middle vertices. Middle edges and neutral vertices, when present,
are exactly those in Proposition~\ref{MM}.

For a side vertex $v$, put
\[
 \delta(v)=d_F(v)+|N(v)\cap B_S|.
\]
Its required number of endpoint-row incidences is $5-\delta(v)$ and
must be positive. Five-vertex sides require no endpoint-row search:
take exactly $\delta(v)$ rows omitting $v$. The sum of these
multiplicities is five. If there is an internal side edge, the tables
force every omission to be one of its endpoints, so these rows are
independent as required.

On a six-vertex side each endpoint row omits a pair. Choose a multiset
of five loopless pairs with degree sequence $\delta$. Every pair must
cover every edge of $D_S$, so its complementary four-subset is
independent. Conversely this condition is sufficient. Parallel pairs
are allowed and represent distinct endpoint rows. This is the same
omission-multigraph leg used in the retained bipartite and singleton
constructions; its unrestricted parameter domain has
$\binom{15+5-1}{5}=11\,628$ multisets. It is not a new large completion
search, nor is $11\,628$ a number of graph classes.

Attach five distinct endpoint-row vertices to the chosen four-subsets
and to $p_0$ or $p_8$. Attach each of the five ports also to $p_4$.
All vertices and repeated equal rows remain distinct.

\begin{lemma}\label{sixradius}
Consider a six-vertex side without middle neighbors, whose internal graph
is one of $K_2,P_3,2K_2,P_4,C_4$, padded with isolates, and whose port
arities are at least two as in the tables. In every feasible incidence
choice, all its side and endpoint-row vertices have distance at most
four from $p_4$.
\end{lemma}
\begin{proof}
Let $Z$ be the side vertices meeting a port. They have distance two
from $p_4$. The endpoint degrees of every internal edge in the listed
graphs have sum at most four. If neither endpoint lay in $Z$, their
endpoint-row incidence sum would be at least six, contradicting the
five independent rows. Thus $Z$ is a vertex cover. A nonisolated vertex
outside $Z$ has distance at most three from $p_4$.

An endpoint row which meets $Z$ has distance at most three. A row
avoiding $Z$ but meeting a nonisolated side vertex has distance at most
four. The only way to have four isolated side vertices is $D_S=K_2$.
If a row avoiding $Z$ consisted of all four of them, all four would
have zero port incidence and hence endpoint incidence five. Every port would
then lie in the two edge endpoints, contrary to independence and its
arity at least two. Hence every endpoint row has distance at most four.
A side vertex isolated in $D_S$ and outside $Z$ has all five endpoint
rows as neighbors. At least one of those rows meets a port-neighbor
vertex, because that vertex has positive endpoint incidence. It follows
that this isolated side vertex has distance at most four as well.
\end{proof}

\begin{lemma}\label{fiveradius}
Every feasible five-vertex side in the tables has all its vertices and
endpoint rows within distance four of $p_4$.
\end{lemma}
\begin{proof}
Every middle vertex has a port neighbor, since its $F$ degree is at
most four. It is therefore at distance two from $p_4$. If the side has
no internal edge, let $Z$ be its vertices meeting a port or middle
vertex. Their distances are at most three. The sum of the corresponding
column deficits is five, and positivity bounds each deficit by four;
thus $|Z|\geq2$. Every endpoint four-subset meets $Z$, so lies within
distance four. A side vertex outside $Z$ meets all five endpoint rows,
including a row incident with a side-port neighbor; its distance is
at most four.

If there is an internal edge, only its two endpoints meet ports or
middle vertices. Each has distance at most three unless its external
deficit is zero. In that event the other endpoint has a port neighbor
and distance two, giving distance three for the zero-deficit endpoint.
Every independent endpoint four-subset contains one of the two endpoints,
so again lies within distance four. The three isolated side vertices
meet every endpoint row, including a row meeting a port-neighbor
endpoint, and have distance at most four.
\end{proof}

For the six-vertex side containing the nonisolated middle vertex in the
$q=2$ table, use the same external support $Z$. In its degree-two row,
the two middle-neighbor vertices and at least two port-neighbor vertices
are disjoint, so $|Z|\geq4$. In its degree-three row the column deficit
sum is ten and each deficit is at most four, so $|Z|\geq3$.
Every endpoint four-subset meets $Z$. Any side edge has at least one
endpoint in $Z$, since otherwise the two endpoint-row incidences sum
to eight. Vertices isolated from both $Z$ and the side graph meet all
five endpoint rows and hence a row at distance three through a port.
This proves the same radius-four conclusion for this last side type.

\begin{theorem}\label{construction}
The construction above, with all its incidence and positivity tests,
gives precisely all nonbipartite cases with a nonisolated middle vertex
relative to the fixed eight-geodesic. Every constructed graph is simple,
quintic, connected, of girth four and diameter eight. No separate global
diameter test is needed. Up to endpoint reversal there are twenty-four
scalar core rows: two for $q=2$, seven for $q=3$, and fifteen for $q=4$.
The rows are not asserted to be disjoint ordinary graph classes.
\end{theorem}
\begin{proof}
Necessity is established by the preceding reductions and case lists.
Conversely each named vertex and each incidence-row vertex has degree
five. Middle degrees are prescribed, and side degrees are five by the
column equations. There is no triangle inside $F$ or in a row's
neighbor subset; these are all the possible triangles. Two endpoint
four-subsets share a side vertex because the sides have order five or
six. Together with their endpoint and two distinct row vertices this
gives a four-cycle. Thus the girth is exactly four.

All middle vertices have a port neighbor, and all port vertices meet
$p_4$. Lemmas~\ref{sixradius} and~\ref{fiveradius}, together with the
last six-vertex argument, put every vertex within four of $p_4$.
Connectivity follows and the diameter is at most eight. The separation
conditions in Lemma~\ref{distance} make the endpoint distance eight,
so the diameter is exactly eight. None of the listed residual graphs
is a five-edge star, hence Lemma~\ref{star} proves nonbipartiteness.
The scalar row counts are $2$, $5+2$, and $7+3+5$, respectively.
\end{proof}

\section{Feasibility and residual shapes}

For completeness we indicate why the tables contain actual incidence
choices, rather than merely integer possibilities. On five-vertex sides
the endpoint rows are forced by the deficits. In every $q=4$ row choose
the middle-star neighbors with their prescribed degrees. A port adjacent
to all middle vertices on its side chooses its side neighbors in the
complement of their union. The displayed arities fit: the only unions
of two degree-two neighborhoods have size at most four, and the unions
of three degree-one neighborhoods have size at most three. All other
port choices are less restrictive. The largest possible column deficit
is four. For the internal-edge rows the incidences at its endpoints
specified in Proposition~\ref{q4} give the required positive columns.
Thus every row has a feasible choice.

Here are small explicit six-vertex legs. For one internal edge $xy$ on
$\{x,y,a,b,c,d\}$ and four ports of arity two, use port subsets
$xa,xb,yc,yd$. Use endpoint rows $abcd$, $xbcd$, $xacd$, $yabd$, $yabc$.
For $P_3$ with centre $v$, leaves $s,t$ and isolates $a,b,c$, and three
ports of arity two, use $va,sb,tc$. The endpoint rows are two copies
of $vabc$ and the three rows consisting of $s,t$ and two of $a,b,c$.
For two disjoint edges $ab,cd$ and isolates $x,y$, use port pairs
$ac,ad,bc$ and endpoint pairs $ac$ twice and $bd$ three times, each
endpoint pair completed by $x,y$.

For $P_4$ in order $0,1,2,3$, all rows contain its two isolates and
an independent path pair. Across the five endpoint rows and two port
rows, the multiplicities of $02,03,13$ must be $(3,1,3)$.
Take both ports to be $02$ and the complementary five endpoint rows.
For $C_4$ and two isolates, take the port to be one opposite pair,
with two endpoint rows using that pair and three using the other pair,
always including both isolates.

The two-port, three-subset $P_3$ leg in the $q=2$ table is exactly the
six-form leg retained from~\cite{middle}. For its other side when the
middle degree is three, take star neighbors $u,v,w$, other vertices
$x,y,z$, the two middle-adjacent port subsets $xy,yz$, and the remaining
port subset $uvw$. Omit from the endpoint rows the four edges of the
cycle $(u,v,w,y)$ and the edge $xz$. For middle degree two, let its
neighbors be $p,q$. If the internal edge is $px$, use ports $xy,xz,yw$;
take endpoint rows $pqzw$ twice, $pyzw$, $xqyw$, $xqyz$.
If the internal edge is disjoint, say $xy$, use ports $xz,xw,yz$;
take endpoint rows $xqzw$, $xpzw$, $ypqw$ twice, $ypqz$.
These choices exhibit both permitted placements of that edge.

The shape list can also be read directly from the tables. The two
$q=2$ rows give $K_{1,3}\sq P_3$, or $P_4\sq P_3$ and
$2P_3\sq K_2$ according as its extra side edge meets a star leaf or
is disjoint. The $q=3$ table gives $K_{1,4}\sq K_2$,
$K_{1,3}\sq P_3$, $K_{1,3}\sq2K_2$, $P_3\sq P_4$ and
$K_2\sq C_4$. Every middle-edge row gives $K_{1,4}\sq K_2$.
In the $q=4$ table, $Q_1,Q_2,Q_3$ give the first two of these shapes;
$Q_4$ gives $K_{1,3}\sq P_3$ or $K_{1,3}\sq2K_2$.
For $Q_5$, the two degree-two side neighborhoods meet in zero, one,
or two vertices, giving $2P_3$, $P_5$, or $C_4$ on that side.
For $Q_6$ the degree-one neighborhood either meets or avoids the
degree-two neighborhood, giving $P_4$ or $P_3\sq K_2$.
For $Q_7$ its three degree-one neighbors have one, two, or three
distinct values, giving $K_{1,3}$, $P_3\sq K_2$, or $3K_2$.
The three internal-edge rows give $K_{1,4}\sq K_2$ or
$K_{1,3}\sq P_3$. Thus the complete possible nonisolated residual shapes, padded
with isolates, are exactly
\[
\begin{gathered}
 K_{1,4}\sq K_2,\quad K_{1,3}\sq P_3,\quad
 K_{1,3}\sq2K_2,\quad P_4\sq P_3,\\
 2P_3\sq K_2,\quad C_4\sq K_2,\quad
 P_5\sq K_2,\quad P_3\sq3K_2.
\end{gathered}
\]
For the last two, use $Q_5$ with the two degree-two neighborhoods
sharing exactly one vertex, or $Q_7$ with its three degree-one neighbors
distinct. The other overlap choices give the preceding shapes. These
eight shapes are a residual-graph list, not eight ordinary graph classes.

\section{A centre which is not a twin class}

Row $Q_4$ makes the distinction particularly clear. Write its middle
vertices as $u,v,w,h$, with $h$ isolated in $F$, $u$ joined to three
vertices on the left, and $v,w$ each joined to one vertex on the right.
Let $B_L=\{a\}$, $B_R=\{b_1,b_2,b_3\}$ and $B_0=\{b_0\}$.
The middle--port table says that $u$ meets $a,b_0$, each of $v,w$
meets $b_0,b_1,b_2,b_3$, and $h$ meets all five ports.
The port $a$ meets the two left vertices outside $N_F(u)$, and the
other three side ports meet singleton right vertices avoiding
$N_F(v)\cup N_F(w)$. The two right middle neighbors may coincide.
Left endpoint rows omit each left vertex once. Right endpoint rows
omit a vertex with multiplicity equal to its middle-plus-port incidence;
if the two middle neighbors coincide at $x$, this omits $x$ twice and
distributes three further omissions according to the three side ports.

The construction theorem gives girth four, nonbipartiteness and diameter
eight. Its centre is exactly $\{p_4,b_0,h\}$ and induces a path of length
two. Indeed $p_4,h$ are false twins. The vertex $b_0$ has distance at
most four from everything: its middle neighbors give the left star and
the right middle neighbors distance two; its distance to each endpoint
is four. A right endpoint row either meets a right middle neighbor or,
if it misses the sole shared one, contains every other side vertex and
therefore a side-port neighbor. A right vertex with neither middle nor
port incidence meets all five endpoint rows, including one of the former
kind. These observations give distances at most four for the remaining
vertices; the left side is immediate from the three-leaf star and its
two-vertex complement.

Every other vertex is within distance at most three of at least one
endpoint: this is immediate for endpoint rows, side vertices and
side ports, and follows through its star edge for $u,v,w$. It therefore
has distance at least five from the opposite endpoint. Only the three
displayed vertices are central. The edges among them are
$p_4b_0,b_0h$, with no edge $p_4h$. In particular a canonicalization
or gluing argument may not replace the entire centre by one false-twin
class.

\paragraph{Remaining computational duty.}
The finite incidence construction still requires actual exhaustive
enumeration, complete ordinary uncolored identification across all roots
and rows, and comparison with the other retained strata before it can
finish the original census. No value of that census, no new case payment,
and no conjecture or formal-verification credit is asserted here.

\begin{thebibliography}{9}
\bibitem{middle} Human Compression project,
\emph{The empty and singleton middle cases for quintic graphs}, frozen
source SHA256 \texttt{76a5c09c2e3f52536c022572bd95ba611392ad38527c0f2ee40c1e886e9d9eb2},
with its separate independent mathematical review.
\bibitem{cages} Cambie, Goedgebeur, Jooken and Van den Eede,
\emph{On the order-diameter ratio of girth-diameter cages},
arXiv:2511.21144v3, Appendix B, Table 6.
\end{thebibliography}
\end{document}
